\documentclass[10pt,a4paper]{report}

\usepackage[utf8]{inputenc}
\usepackage[francais]{babel}
\usepackage[T1]{fontenc}
\usepackage{lipsum}
\usepackage{fancyhdr}
\usepackage{listings}
\usepackage{graphicx}
\usepackage{geometry}
\usepackage{hyperref}
\usepackage[dvipsnames]{xcolor}
\usepackage{setspace}

\hypersetup{
	colorlinks,
	urlcolor=black
	filecolor=black,
	linkcolor=black,
	citecolor=black,
}

\definecolor{mygray}{RGB}{233,233,233}

\lstset{
	language=Prolog,
	basicstyle=\ttfamily\small,
	identifierstyle=\color{Cerulean},
	numberstyle=\color{green},
	keywordstyle=\color{blue},
	commentstyle=\it\color{green},
	stringstyle=\color{black!60},
	columns=flexible,
	tabsize=4,
	extendedchars=true,
	showspaces=false,
	showstringspaces=false,
	numbers=left,
	numberstyle=\small,
	breaklines=true,
	breakautoindent=true,
	captionpos=b,
	backgroundcolor=\color{mygray},
	captionpos = b 
}

\newcommand{\HRule}{\rule{\linewidth}{1mm}}
\newsavebox{\BBbox}
\newenvironment{DDbox}[1]{
	\begin{lrbox}{\BBbox}\begin{minipage}{\linewidth}}
		{\end{minipage}\end{lrbox}\noindent\colorbox{mygray}{\usebox{\BBbox}} \\
	[.5cm]}

\fancypagestyle{plain}{
	\fancyhf{}
	\fancyhead[L]{Soutenance}
	\fancyhead[R]{\ifnum\value{part}>0 \partname\ \thepart \fi}
	\renewcommand{\headrulewidth}{0.4pt}
	\fancyfoot[L]{Aurélien Thirion et Nicolas Blin}
	\fancyfoot[C]{\thepage}
	\fancyfoot[R]{\today}
	\renewcommand{\footrulewidth}{0.4pt}
}


\geometry{
	a4paper,
	total={210mm,297mm},
	left=20mm,
	right=20mm,
	top=20mm,
	bottom=20mm,
}

\begin{document}
\begin{titlepage}
    \begin{center}
      ~\\[4cm]
      \textsc{\huge Université de Lorraine}\\[1.0cm]
      \includegraphics{univr.jpg}\\[0.5cm]
      \textsc{\Large Faculté de Science et Technologie}\\[0.5cm]
      \textsc{\Large Master 1 Informatique}\\[0.5cm]
      \HRule \\[0.4cm]
      { \huge \bfseries Projet de LMC \\[0.4cm] }
      \HRule \\[1cm]
     \textsc{\LARGE Algorithme d'unification \\ ~ \\ Martelli Montanari }\\[4cm]
	\emph{Auteurs:}\\
	Aurélien \textsc{Thirion} \\
	Nicolas \textsc{Blin}\\
      \vfill
      {\large \today}
    \end{center}
  \end{titlepage}








\chapter*{Introduction}
Nous avons voulu, dans ce projet, essayer d'être le plus clair et simple possible en
disposant dans différents fichiers le code source de notre programme. Ainsi dans le code
principal, il y a le minimum de ligne pour comprendre réellement de processus de
l'unification de Martelli-Montanari. Notre code principal se situe dans \textbf{main.pl}, des
prédicats sont appelés lors du processus dans les fichiers suivants : \textbf{operateur.pl} ,
\textbf{predicatsRelais.pl} et \textbf{reglesTest.pl}.\\[1cm]

\begin{center}
\textbf{operateur.pl}\\
Contient l'opérateur \textbf{?=} et le code pour le echo contenu dans le mail.\\[1cm]
\textbf{predicatsRelais.pl}\\
Fonction de décomposition d'une équation E pour obtenir la partie gauche de
l'opérateur \textbf{?=} et la partie droite.\\[1cm]
\textbf{reglesTest.pl}\\
Toutes les conditions nécessaire pour appliquer la réduction (prédicat reduit(R,E,P,Q)).\\[1cm]
\end{center}





\chapter*{Question 1}
\paragraph{Mise en place pour l'unification} ~\\
On peut voir ici les instructions pour l'ouverture des différents fichiers.
\begin{lstlisting}[caption ={Dans le fichier $main.pl$}]
:-
	[operateurs],
	[predicatsRelais],
	[reglesTest].
\end{lstlisting} ~\\
Un exemple du prédicat regle(E,R) dans le code.
\begin{lstlisting}[caption ={Dans le fichier $main.pl$}]
regle(E,rename)
\end{lstlisting} ~\\
Qui va appeler le prédicat dans le fichier \textbf{reglesTest.pl}. Ici on va séparer l'équation E en les termes X et T. Il va ensuite vérifier l'applicabilité de la règle $rename$.
\begin{lstlisting}[caption ={Dans le fichier $reglesTest.pl$}]
regle(E, rename):-
	splitEquation(E,X,T),
	var(T),
	var(X).
\end{lstlisting} ~\\
Qui va appeler le prédicat dans le fichier \textbf{predicatsRelais.pl}. On récupère la partie gauche et droite de \textbf{?=} pour le mettre dans X et T respectivement.
\begin{lstlisting}[caption ={Dans le fichier $predicatsRelais.pl$}]
splitEquation(E,X,T):-
	arg(1,E,L),
	arg(2,E,R),
	X = L,
	T = R.
\end{lstlisting}

\paragraph{Occur check} ~\\
Le prédicat $occur_check$ a été assez simple à mettre en oeuvre. Il sera chargé de vérifier si V apparaît dans le terme composé T.
\begin{lstlisting}[caption ={Dans le fichier $main.pl$}]
occur_check(V,T):-
	compound(T),
	var(V),
	contains_var(V,T).
\end{lstlisting}

\paragraph{Reduit} ~\\
Voici un exemple de $reduit$ qui se charge d'appliquer la règle rename.
\begin{lstlisting}[caption ={reduit dans $main.pl$}]
reduit(rename, E, P, Q):-
	splitEquation(E,X,T),
	X = T,
	P = [_|Q].
\end{lstlisting}





\chapter*{Question 2}
\paragraph{Dans le main} ~\\





\mbox{}
\newpage
\chapter*{Code source}
\paragraph{Dans le main} ~\\
\begin{lstlisting}[caption ={$main.pl$}]
occur_check(V,T):-
	compound(T),
	var(V),
	contains_var(V,T).

unif(P,S) :-
	clr_echo,
	unifie(P,S).

trace_unif(P,S) :-
	set_echo,
	(unifie(P,S),
	 echo("Yes"),
	 !;
	 echo("No")).

unifie([], _) :- !.
unifie([]) :- !.


unifie(P):-
	unifie(P, regle),
	!.

unifie(P, regle):- unifie(P, rename).
unifie(P, regle):- unifie(P, simplify).
unifie(P, regle):- unifie(P, expand).
unifie(P, regle):- unifie(P, check).
unifie(P, regle):- unifie(P, orient).
unifie(P, regle):- unifie(P, decompose).
unifie(P, regle):- unifie(P, clash).


unifie(P, rename) :- 
	P = [E |_],
	regle(E, rename),
	reduit(rename, E, P, Q),
	unifie(Q, regle),!.


unifie(P, simplify):-
	P = [E |_],
	regle(E, simplify),
	reduit(simplify, E, P, Q),
	unifie(Q, regle),!.

unifie(P, expand):-
	P = [E |_],
	regle(E, expand),
	reduit(expand, E, P, Q),
	unifie(Q, regle),!.

unifie(P, check):-
	P = [E |_],
	regle(E, check),
	reduit(check, E, P, Q),
	unifie(Q, regle),!.

unifie(P, orient):-
	P = [E |_],
	regle(E, orient),
	reduit(orient, E, P, Q),
	unifie(Q, regle),!.

unifie(P, decompose):-
	P = [E |_],
	regle(E, decompose),
	reduit(decompose, E, P, Q),
	unifie(Q, regle),!.

unifie(P, clash):-
	P = [E |_],
	regle(E, clash),
	reduit(clash, E, P, Q),
	unifie(Q, regle),!.


reduit(decompose, E, P, Q):-
	splitEquation(E,X,T),
	functor(X,_,ArityX),
	functor(T,_,_),
	P = [_|Tail],
	repet(X,T,ArityX,Tail,Q),
	echo("decompose: "),echo(Q),nl.

repet(_,_,0,T,Q):- Q = T, !.
repet(X,T,N,Tail,Q) :-
	N > 0,
	arg(N,X,ValX),
	arg(N,T,ValT),
	Var = [ValX?=ValT|Tail],
	N1 is N - 1,
	repet(X,T,N1,Var,Q).

reduit(rename, E, P, Q):-
	splitEquation(E,X,T),
	X = T,
	P = [_|Q].

reduit(simplify, E, P, Q):-
	splitEquation(E,X,T),
	X = T,
	P = [_|Q].

reduit(expand, E, P, Q):-
	splitEquation(E,X,T),
	X = T,
	P = [_|Q].

reduit(check, _, _, _):-
	fail,
	!.

reduit(orient, E, P, Q):-
	splitEquation(E,X,T),
	P = [_|Tail],
	Q = [T ?= X | Tail].


reduit(clash, _, _, _):-
	fail,
	!.
\end{lstlisting}



\end{document}
